\section{Training and Evaluation Details}
\label{app:details}

\paragraph{Hyperparameters.}
All PhyWorld runs: AdamW, lr $5{\times}10^{-5}$ (from-scratch reruns: $2{\times}10^{-5}$), weight decay $10^{-3}$, bf16, gradient clip $1.0$, $20$ epochs ($60$ for from-scratch cells; the pretraining-regime $2{\times}2$ uses $60$ epochs uniformly), batch size $64$ for free rollout ($H{=}8$), $128$ for teacher forcing ($H{=}1$), $48$ for $H{\geq}16$. History size $3$; embedding dimension $192$; SIGReg weight $0.09$. Seeds: $3072/1234/42$ where seed-replicated. Physion++ runs use the readout split ($800$ clips) with the same optimizer.

\paragraph{OOD partitions.}
ID: radius/mass $\in[0.7,1.5]$, velocity $\in[1.0,4.0]$. OOD draws radius/mass from $[0.3,0.6]\cup[1.5,2.0]$ and velocity from $[0,0.8]\cup[4.5,6.0]$. Collision uses the four-column variant (two objects' masses and velocities). Training sets are ID-only ($1{,}000$ trajectories); evaluation covers all four partitions ($\geq$1{,}600 trajectories per domain).

\paragraph{Probe protocol.}
Ridge regression ($\alpha{=}1$), trajectory-level 80/20 splits, probing the encoder output (not the projector), pretrained ``paper'' initialization, and $K{=}4$ stacked consecutive latents for velocity. Each choice matters: the naive combination (random init, single frame, through-projector) reads $\rho{=}0.166$ for uniform velocity where the corrected protocol reads $0.939$ (Trap 4, Appendix~\ref{app:traps}). Random-encoder and pixel-statistics controls are run alongside every probe.

\section{Full Injection Results}
\label{app:injection}

\paragraph{Probe/slot factorial (uniform, free rollout).}
Latent both-OOD nMSE / pixel both-OOD PSNR (dB): baseline $0.131$ / $20.41$; probe-only $0.167$ / $19.83$; slot$+$reweight $0.114$ / $21.30$; probe$+$slot$+$reweight $0.125$ / $20.02$. The combination underperforms the slot alone on both scales; the latent and pixel verdicts agree.

\paragraph{Slot reweighting sweep (uniform).}
Slot weight $1/30/100$: $0.183$ / $0.114$ / $0.136$ (single seed), with the $w{=}30$ cell seed-replicated at $0.132\pm0.017$ against the baseline $0.136\pm0.008$. Kinematic head on the reweighted slot: velocity-OOD decoded position $\rho{=}0.991/0.987$ ($x$/$y$); uniform pixel horizon-28 $23.34$ vs.\ $22.09$\,dB ($+1.25$); appearance-OOD remains the weak axis (decoded $\rho$ $0.29\!\to\!0.43$, vs.\ $0.96$ without the kinematic head).

\paragraph{Bypass probe details.}
Decoded position $\rho$ per domain (uniform/parabola/collision). Baseline model: full $192$ dims $0.92/0.85/0.78$; black-box $[2{:}192]$ alone $0.92/0.85/0.79$. Slot-pinned model (structpos, $w{=}30$): slot dims $[0{:}2]$ alone $0.92/0.85/0.51$; black-box $[2{:}192]$ alone $0.92/0.86/0.79$---unchanged by pinning. The slot absorbs position (its supervision loss falls $2.53\!\to\!0.015$) without displacing the black box's redundant copy.

\paragraph{Encoded quantity vs.\ domain (position monotone, velocity sign-flipping).}
The reweighted slot's effect is systematic in what it encodes (hardest clean split; baseline: uniform both-OOD $0.131$, parabola \rmood{} $0.127$, collision both-OOD $0.393$):

\begin{center}\small
\begin{tabular}{@{}lccc@{}}
\toprule
slot content & uniform & parabola & collision \\
\midrule
position (structpos) & $0.132$ & $0.160$ & $0.596$ \\
position$+$velocity & $0.207$ & $0.096$ & $0.621$ \\
$\Delta$ (adding velocity) & $+0.075$ & $-0.064$ & $+0.025$ \\
\bottomrule
\end{tabular}
\end{center}

Position is an output variable with the same role in every domain, so pinning it degrades monotonically with dynamics complexity and never flips sign. The marginal effect of adding velocity ($+0.075/{-}0.064/{+}0.025$) tracks velocity's role in each domain---redundant constant (uniform), linear driver (parabola), discontinuous jump (collision). The parabola gain ($0.122\pm0.007\!\to\!0.096\pm0.007$, every seed lower) is thus the sole seed-robust structural gain in the grid: a fixed form landing on a matching domain, not a portable prior.

\paragraph{Physion++ per-scenario (all ten scenarios, FR at $20$ epochs).}
Rollout nMSE: bouncy\_platform $0.041$, bouncy\_wall $0.094$, friction\_cloth $0.071$, friction\_collision $0.062$, friction\_platform $0.041$, mass\_collision $0.083$, mass\_dominoes $0.058$, mass\_waterpush $0.122$; deformable scenes are the common hard case for every method (deform\_clothhang $1.375$, deform\_clothhit $0.772$). Injection arms degrade the rigid-body scenarios by $3$--$10\times$ (Table~\ref{tab:pp}) while leaving the deformable failure in place.

\paragraph{Physion++ horizons.}
Overall nMSE at horizons $16/32/64$: baseline ($H{=}8$) $0.016/0.140/0.280$; $H{=}20$ $0.022/0.033/0.220$; $H{=}16{+}$scale $0.010/0.035/0.136$; $H{=}20{+}$scale $0.012/0.024/0.087$; $H{=}28{+}$scale reaches horizon-$64$ nMSE $0.014$. Free rollout is the main long-horizon lever (horizon-$32$ cosine $0.91$ vs.\ $0.79$--$0.87$ for every physics-augmented variant); geometric scale jitter acts as an amplitude stabilizer for long rollouts: without it, $H{=}20$ explodes on friction-collision (nMSE $24.6$ at cosine $0.99$); with it, $0.019$.

\section{Augmentation: a Domain-Specific Lever with a Reversal Boundary}
\label{app:aug}

Augmentation is not part of the paper's main line (it is neither an injection mechanism nor a protocol variable), but its boundary behavior is a useful caution. On synthetic domains, per-trial appearance jitter (brightness/contrast, strength $0.5$) roughly halves the hardest split: uniform $0.131\!\to\!0.068$ ($-48\%$), parabola $-63\%$ (\rmood{} $0.127\!\to\!0.065$); geometric scale jitter, only in combination with a long horizon, sets the collision record ($0.393\!\to\!0.172\pm0.004$ over three seeds). Interactions are non-monotone: appearance conflicts with long-horizon training ($0.472$, worse than either alone), scale synergizes ($0.208$), the triple is worse than the best pair ($0.253$); temporal-stride augmentation loses even to appearance on velocity-OOD ($0.556$ vs.\ $0.376$), leaving unseen velocities the open hard case.

The appearance recipe reverses on photorealistic simulation: the same jitter that halves synthetic OOD error degrades Physion++ friction-collision prediction by two orders of magnitude (nMSE $0.062\!\to\!6.44$) while the \emph{aggregate} long-horizon cosine improves (Trap 1 below), and it does not help zero-shot transfer either ($0.597$, below the $0.607$ random-prior ceiling). The boundary is principled: appearance jitter encodes an appearance-invariance assumption that holds under PhyWorld's simple renderer but breaks where appearance carries physical information (friction, mass, material). Scale jitter, which assumes only size-invariance, keeps helping at photorealistic scale (Appendix~\ref{app:injection}); appearance jitter flips sign. Augmentation gains are domain- and type-specific, not portable.

\section{Metric Pathologies and Evaluation Traps}
\label{app:traps}

Our anatomy repeatedly reached conclusions opposite to those suggested by common metrics; the decision rules of \S\ref{sec:setup} distill the following four failure modes, each of which reversed at least one finding.

\paragraph{Trap 1: cosine and probe metrics are duals of the training loss.}
Probe correlation and latent cosine rise mechanically when probe or slot losses are optimized---they measure whether information is \emph{present}, not whether prediction \emph{uses} it. Probe supervision lifts velocity decodability from $\rho{=}0.44$ to $0.80$ and gives the most stable long-horizon \emph{cosine} ($0.882$ vs.\ $0.843$), while its \emph{pixel} rollout is the worst in the comparison ($19.93$ vs.\ $20.64$\,dB at horizon $28$). Appearance augmentation on Physion++ improves the aggregate long-horizon cosine while per-scene nMSE collapses up to $100\times$ (deform\_clothhit: cosine $0.610\!\to\!0.913$, nMSE $0.772\!\to\!3.69$). Judging either result by direction-only or readout metrics inverts the conclusion. Our own early sweep was misled this way: judged on $K{=}4$ probe $\rho$, $\lambda{=}50$ ``won''; re-judged on prediction loss, the ranking reversed.

\paragraph{Trap 2: nMSE denominators can degenerate.}
When the target sequence degenerates (a parabola ball exits the frame, target variance ${\to}0$), per-horizon nMSE explodes---$3{\times}10^{4}$--$2{\times}10^{6}$ at horizon ${\sim}28$ across all six parabola baseline arms, while cosine sits at a healthy $0.55$--$0.95$---and pollutes every aggregate that includes late horizons. This is why parabola is judged on its clean \rmood{} partition throughout. Traps 1 and 2 point in opposite directions (cosine under-reports failure; degenerate nMSE over-reports it); neither metric family is safe alone, hence per-partition, per-horizon cross-validation.

\paragraph{Trap 3: zero-shot transfer is bounded by the architecture prior.}
On Physion OCP, a \emph{randomly initialized} encoder scores mean AUC $0.607$. No synthetic-training configuration exceeds it: free rollout approaches it ($0.603$), physics variants fall below it ($0.551$--$0.582$), appearance augmentation $0.579$--$0.597$; training longer approaches the ceiling monotonically ($0.554$ at epoch $1$ to $0.603$ at epoch $20$) without crossing. Per-scenario, only Support ($+0.10$) and Collide ($+0.07$) show signal above the random control. Zero-shot claims must be calibrated against the random-prior floor, or they measure the architecture, not the learning.

\paragraph{Trap 4: protocol confounds fabricate results in both directions.}
Three probe-protocol choices (random vs.\ pretrained initialization, single-frame vs.\ stacked probes, probing through the projector) multiply into a false negative: the naive stack reads uniform velocity decodability at $R^{2}{=}0.166$, the corrected protocol $0.939$. A silent initialization bug (only $24$ of $216$ loadable encoder keys actually loaded) invalidated an entire 45-configuration sweep before a loading guard caught it. And converged training loss proves nothing about rollout: our clean from-scratch baseline converges to the pretrained model's prediction loss ($0.006$ vs.\ $0.005$) while its rollout OOD error is $2.8\times$ worse.

\paragraph{Auxiliary-loss dominance (measurement behind \S\ref{sec:mech}, Test 2).}
At probe weight $\lambda{=}10$, the weighted auxiliary-to-prediction loss ratios---used as a proxy for gradient magnitude; we do not measure gradient norms directly---are $15\times$ (collision), $44\times$ (uniform), and $125\times$ (parabola); the encoder's effective (participation-ratio) dimensionality collapses by $39$--$90\%$ relative to baseline (uniform $41\!\to\!4$). At $\lambda{=}1$ the ratios are benign, which is the setting used in \S\ref{sec:scan}.

\paragraph{Checklist.}
(i)~Report per-partition \emph{and} per-horizon; never aggregate across horizons without inspecting them. (ii)~Judge with scale-aware prediction metrics (nMSE, decoded-pixel PSNR); use cosine and probe $\rho$ as diagnostics only. (iii)~Audit nMSE denominators for degenerate targets. (iv)~Calibrate transfer against a random-init control, and validate probe protocols on a random encoder.
